\markdownRendererDocumentBegin
\markdownRendererSectionBegin
\markdownRendererHeadingOne{Exponenciación Binaria}\markdownRendererInterblockSeparator
{}La operación \markdownRendererDollarSign{}a\markdownRendererCircumflex{}b\markdownRendererDollarSign{} en la gran mayoría de librerías se implementa en \markdownRendererDollarSign{}O(b)\markdownRendererDollarSign{}.\markdownRendererHardLineBreak
{}Sin embargo, si representamos \markdownRendererDollarSign{}b\markdownRendererDollarSign{} en base 2 (binario), se puede realizar en \markdownRendererDollarSign{}O(log b)\markdownRendererDollarSign{}:\markdownRendererInterblockSeparator
{}\markdownRendererInputFencedCode{./_markdown_main/745e27c31ab3cb2b79f8ec72c1a6ff98.verbatim}{cpp}{cpp}\markdownRendererInterblockSeparator
{}
\markdownRendererSectionEnd \markdownRendererSectionBegin
\markdownRendererHeadingOne{Aritmética Modular}\markdownRendererInterblockSeparator
{}Cuando los números de entrada son muy grandes, los problemas suelen requerir que apliquemos aritmética modular con un número p, típicamente \markdownRendererDollarSign{}10\markdownRendererCircumflex{}9\markdownRendererDollarSign{} + 7 (1000000007).\markdownRendererParagraphSeparator
{}Por ejemplo, supongamos que queremos calcular: \markdownRendererDollarSign{}r = (a * b * c) \markdownRendererPercentSign{} 1000000007;\markdownRendererDollarSign{}\markdownRendererParagraphSeparator
{}Si realizamos la multiplicación directamente, los números intermedios pueden exceder la capacidad de almacenamiento, aunque matemáticamente sea correcto. Por eso usamos operaciones modulares paso a paso.\markdownRendererInterblockSeparator
{}\markdownRendererSectionBegin
\markdownRendererSectionBegin
\markdownRendererHeadingThree{Operaciones mediante aritmética modular}\markdownRendererParagraphSeparator
{}\markdownRendererInputFencedCode{./_markdown_main/810baa4c640570705ab22502c127fc11.verbatim}{cpp}{cpp}
\markdownRendererSectionEnd 
\markdownRendererSectionEnd 
\markdownRendererSectionEnd \markdownRendererDocumentEnd